\section{Conclusion}
\label{sec:conclusion}

\begingroup\color{blue}
Promotion frequency and arrangement jointly determine revenue because historical demand effects interact with transaction prices. The window contribution functions translate these interactions into sufficient conditions for balanced or single block optimality at prescribed cycle length and promotion count. They also identify when frequency choice admits a finite search or reduces to comparing constant prices.
\par\endgroup

\begingroup\color{blue}
At prescribed cycle length and promotion count, arrangement optimization is NP-hard in general under mixed curvature, while short-cycle searches can miss the unrestricted optimum. Finite searches over balanced and single block policies, together with memory truncation, provide alternatives with revenue-loss guarantees. In the tested synthetic instances, simple policies perform well despite conservative bounds, while the effectiveness of truncation depends on how strongly older histories affect demand. The access-service example also shows how arrangement can change revenue and demand variability at the same frequency.
\par\endgroup

\begingroup\color{blue}
The model takes customer arrivals and the two prices as given and excludes strategic purchase delay and operating constraints. Incorporating operating costs and capacity constraints would help quantify the value of the smoother demand profiles observed under balanced arrangements. Accounting for demand estimation uncertainty and using field data would further assess the structural conditions and the informativeness of the revenue-loss guarantees. The complexity of unrestricted frequency and arrangement optimization remains open.
\par\endgroup
